\section{Retrieval System}

\label{sec:retrieval}

\subsection{Embedding Space}

Each experience $e \in \mathcal{E}$ is embedded into $\mathbb{R}^{7680}$ via the Zen-Reranker model~\cite{zoo2025zen}:

\begin{equation}
e.\text{emb} = \text{ZenEmbed}(e.\text{text})
\end{equation}

\textbf{Zen-Reranker Properties}:
\begin{itemize}
    \item Architecture: Dense retrieval with cross-attention
    \item Dimension: 7680 (vs. 1536 for OpenAI, 4096 for Nomic)
    \item Training: Multilingual, cross-modal (text/code/math)
    \item Precision: FP16 (15 KB per embedding)
\end{itemize}

\subsection{Similarity Search}

Given query $q$, we retrieve top-$k$ relevant experiences via cosine similarity:

\begin{equation}
\text{sim}(q, e) = \frac{q.\text{emb} \cdot e.\text{emb}}{\|q.\text{emb}\| \|e.\text{emb}\|}
\end{equation}

\begin{algorithm}[t]
\caption{Top-k Experience Retrieval}
\label{alg:retrieval}
\begin{algorithmic}[1]
\STATE \textbf{Input}: Query $q$, library $\mathcal{E}$, $k$
\STATE \textbf{Output}: Top-$k$ experiences $\mathcal{E}_k$
\STATE $q.\text{emb} \leftarrow \text{ZenEmbed}(q)$
\STATE $\text{scores} \leftarrow [\text{sim}(q, e) \mid e \in \mathcal{E}]$
\STATE $\text{indices} \leftarrow \text{argsort}(\text{scores}, \text{descending})$
\RETURN $\{e_{\text{indices}[i]} \mid i < k\}$
\end{algorithmic}
\end{algorithm}

\textbf{Optimization}: For large libraries ($|\mathcal{E}| > 1000$), we use FAISS~\cite{johnson2019billion} with HNSW indexing to reduce complexity from $O(n)$ to $O(\log n)$.

\subsection{Context Injection Protocol}

Retrieved experiences are formatted as a prompt prefix:

\begin{lstlisting}[caption=Context Injection Template]
# Learned Experiences

1. When solving geometry problems with 
   intersections, validate solutions lie within 
   bounded regions, not extensions.

2. For expected extreme statistics in 
   combinatorial problems, use direct enumeration 
   for small sizes.

[... top-k experiences ...]

---

# Problem

{user_query}
\end{lstlisting}

\textbf{Hyperparameters}:
\begin{itemize}
    \item $k = 5$ (empirically optimal, Section~\ref{sec:evaluation})
    \item Similarity threshold: $\text{sim} > 0.6$ (filter irrelevant)
    \item Max context length: 2048 tokens (balance relevance vs. cost)
\end{itemize}

\subsection{Coverage Analysis}

\begin{definition}[Experience Coverage]
The coverage $C(\mathcal{E}, \mathcal{Q})$ of library $\mathcal{E}$ over query set $\mathcal{Q}$ is:

\begin{equation}
C(\mathcal{E}, \mathcal{Q}) = \frac{1}{|\mathcal{Q}|} \sum_{q \in \mathcal{Q}} \mathbb{1}[\max_{e \in \mathcal{E}} \text{sim}(q, e) > \tau]
\end{equation}

where $\tau = 0.6$ is the relevance threshold.
\end{definition}

\begin{theorem}[Coverage Growth]
Given library size $n = |\mathcal{E}|$ and uniform distribution over $d$ domains, expected coverage satisfies:

\begin{equation}
\mathbb{E}[C(\mathcal{E}, \mathcal{Q})] \geq 1 - e^{-n/d}
\end{equation}
\end{theorem}

\begin{proof}
Each experience covers an expected fraction $1/d$ of queries (assuming uniform domain distribution). The probability a query is uncovered after $n$ experiences is $(1 - 1/d)^n \approx e^{-n/d}$ for large $d$. Thus coverage is $1 - e^{-n/d}$.
\end{proof}

\textbf{Implication}: To achieve 95\% coverage over $d=20$ domains requires $n \approx 60$ experiences, consistent with empirical observations (Section~\ref{sec:evaluation}).

